\ExperimentEntry
  {2026-04-20}
  {EXP-BASELINE-001}
  {Ретроспективно зафиксирована текущая исследовательская база проекта. В репозитории собраны baseline-методы \texttt{rand+nn} и \texttt{2opt+greed}, более сильные методы \texttt{grasp} и \texttt{lkh-wrapper-v2}, а также единый пайплайн для генерации инстансов, пакетных прогонов и агрегации результатов.}
  {Основные машинные артефакты: \path{data/results/mtsp_results.csv}, \path{data/results/mtsp_summary.csv}, \path{data/results/mtsp_report.md}. Базовая сетка: $n \in \{20, 50, 100, 200\}$, $m \in \{2, 3, 5\}$, по 10 повторов на каждую конфигурацию.}
  {Проект переведен в состояние воспроизводимого исследования: все ключевые числа теперь получаются не вручную, а через фиксированный экспериментальный пайплайн.}

\ExperimentEntry
  {2026-04-20}
  {EXP-REFERENCE-002}
  {Добавлен и зафиксирован единый внешний reference baseline на базе OR-Tools.}
  {Каноническим источником сравнения объявлен файл \path{data/results/mtsp_reference_results.csv}. Все последующие сравнения ведутся только относительно него, а в отчетах используется строковый формат \texttt{our | OR-Tools | gap | time}.}
  {Появилась внешняя точка отсчета: теперь можно измерять не только внутреннее превосходство наших методов, но и расстояние до сильного reference-решения. При этом OR-Tools интерпретируется как внешний baseline, а не как доказанный оптимум.}

\ExperimentEntry
  {2026-04-20}
  {EXP-BENCHMARK-003}
  {Расширен набор синтетических тестов для mTSP.}
  {Сгенерировано 120 инстансов: четыре размера графа, три значения числа коммивояжеров и 10 повторов на каждую пару $(n, m)$. Результаты агрегированы в \path{data/results/mtsp_reference_summary.csv} и связанных CSV/Markdown-артефактах.}
  {Это дает более устойчивую статистическую картину и позволяет сравнивать алгоритмы уже не на нескольких примерах, а на полноценной выборке.}

\ExperimentEntry
  {2026-04-20}
  {EXP-LKHV2-004}
  {Создана версия \texttt{lkh-wrapper-v2} как отдельное улучшение относительно \texttt{lkh-wrapper-v1}. В версию \texttt{v2} добавлены candidate sets, lookahead при построении начального решения и более направленный локальный поиск.}
  {Сравнение версий ведется через \path{data/results/lkh_versions_reference_summary.csv}. По текущей сводке средний gap уменьшился примерно с $193.803$ до $92.603$, а среднее время снизилось примерно с $0.019135$ до $0.011618$ секунд.}
  {Версия \texttt{v2} оказалась лучше \texttt{v1} на всех 12 конфигурациях, поэтому дальнейшие улучшения имеет смысл строить именно как последовательность новых версий \texttt{v3}, \texttt{v4} и так далее с обязательной фиксацией gap к каноническому reference baseline.}

\ExperimentEntry
  {2026-04-20}
  {EXP-MAINCFG-005}
  {Основной benchmark проекта переведен с \texttt{lkh-wrapper} на \texttt{lkh-wrapper-v2}.}
  {Изменен основной конфиг \path{experiments/config.json}, обновлены поясняющие документы и заново пересчитаны основные CSV- и Markdown-артефакты для \path{data/results/mtsp_*}. Версия \texttt{lkh-wrapper-v1} сохранена только в отдельной линии versioned-сравнения.}
  {После этого основной экспериментальный контур стал согласован с фактическим состоянием исследования: в нем используется не историческая, а актуальная рабочая версия эвристики.}

\ExperimentEntry
  {2026-04-20}
  {EXP-MATH-APPENDIX-006}
  {Подготовлено отдельное LaTeX-приложение с математическим обоснованием именно тех эвристик, которые реально реализованы в репозитории: \texttt{rand+nn}, \texttt{2opt+greed}, \texttt{grasp} и \texttt{lkh-wrapper-v2}.}
  {Новый файл \path{docs/appendices/mtsp_math_appendix.tex} фиксирует постановку задачи, обоснование допустимости маршрутов, монотонность локальных улучшений, условия конечности алгоритмов и оценки сложности, согласованные с текущей реализацией в \path{src/}.}
  {Приложение можно включать в состав курсовой как самостоятельный технический раздел: оно опирается на код проекта и не содержит обзорных фрагментов, не связанных с фактической реализацией.}
